\section{Proposed DecMed-PGHD Module}

\subsection{Architecture and Trust Boundaries}

Fig.~\ref{fig:architecture} shows the extended DecMed architecture.
The red components are introduced or modified for PGHD.
The patient environment contains the data sources, the Android client, and local storage.
The backend contains the PRE service, Redis, IPFS, IOTA smart contracts, and an IOTA Gas Station.
The healthcare environment contains the clinician client.
The PRE service is treated as a semi-trusted coordinator: it may verify and route data, but it must never receive plaintext PGHD or the patient's private key.

\begin{figure*}[t]
	\centering
	\begin{tikzpicture}[
			>=Latex,
			font=\scriptsize,
			component/.style={
					draw,
					rounded corners=2pt,
					align=center,
					minimum height=10mm,
					text width=30mm,
					fill=black!4
				},
			store/.style={
					draw,
					align=center,
					minimum height=8mm,
					text width=22mm,
					fill=black!10
				},
			flow/.style={->,thick},
			control/.style={->,thick,dashed}
		]
		\node[component] (sources) {Wearable/Health Connect\\Device sensors\\Manual input};
		\node[component, right=8mm of sources] (patient) {Patient Android client\\Room + WorkManager};
		\node[component, right=12mm of patient] (pre) {PRE service\\Verification and delegation};
		\node[component, right=12mm of pre] (clinician) {Clinician client\\Verification and decryption};
		\node[store, below=12mm of patient] (iota) {IOTA\\Access, metadata, validity};
		\node[store, below=12mm of pre] (ipfs) {IPFS\\Encrypted PGHD};
		\node[store, below=12mm of clinician] (redis) {Redis\\Temporary access material};

		\draw[flow] (sources) -- (patient);
		\draw[flow] (patient) -- (pre);
		\draw[<->,thick] (pre) -- (clinician);
		\draw[control] (patient) -- node[left,align=right] {grant/revoke} (iota);
		\draw[control] (pre) to[bend right=18] node[above,sloped] {metadata/status} (iota);
		\draw[flow] (pre) -- node[right] {ciphertext} (ipfs);
		\draw[control] (pre) to[bend left=18] node[above,sloped] {temporary material} (redis);
	\end{tikzpicture}
	\caption{End-to-end DecMed-PGHD architecture and data flow. Plaintext PGHD is confined to the patient and authorized clinician endpoints.}
	\label{fig:architecture}
\end{figure*}

The design separates data according to sensitivity, volume, and mutability, as summarized in Table~\ref{tab:data-placement}.
Bulk encrypted content is stored off-chain, while the ledger maintains compact metadata and access state.
Redis stores only temporary access artifacts; it is not an authoritative PGHD repository.

\begin{table}[t]
	\caption{Data Placement in DecMed-PGHD}
	\label{tab:data-placement}
	\centering
	\footnotesize
	\begin{tabular}{@{}p{0.17\columnwidth}p{0.22\columnwidth}p{0.47\columnwidth}@{}}
		\toprule
		\textbf{Location} & \textbf{Representation} & \textbf{Purpose} \\
		\midrule
		Patient device & Records and encrypted batches & Offline retention, batching, retry state \\
		IPFS & Encrypted PGHD payload & Content-addressed bulk storage \\
		IOTA & CID, hash, validity, access, and cryptographic metadata & Audit, authorization, and integrity anchor \\
		Redis & Short-lived token and re-encryption material & Temporary PRE coordination \\
		\bottomrule
	\end{tabular}
\end{table}

The main adversarial conditions considered are an unauthorized clinician, tampering during transit or storage, a semi-trusted infrastructure component attempting to inspect content, temporary network or service failure, and a workload burst.
Compromise of the patient endpoint before encryption and clinical correctness of the source device are outside the present scope.

\subsection{Offline-First Collection and Batching}

The Android client acquires PGHD from three source categories: Health Connect records synchronized from applications or wearables, direct Android sensor readings, and manual patient input.
Each record is immediately persisted in a local Room database with its measurement time, source category, source application or device when available, recording method, value, and unit.
Collection therefore continues independently of network availability.

Batch creation and synchronization are scheduled through WorkManager.
A batch is triggered periodically, after a configured size threshold, manually, or when connectivity returns.
Records are grouped by measurement type and source.
Each group contains raw data points, summary statistics, configured clinical thresholds, and anomaly markers.
The anomaly markers are aids for presentation and do not constitute a diagnosis.

The client maintains explicit \texttt{pending}, \texttt{failed}, and \texttt{sent} states.
Records already assigned to a batch are mapped to that batch so that a retry does not construct a second logical batch from the same records.
Failed work remains in local storage and is rescheduled after the relevant constraint is satisfied.
The transport consequently provides at-least-once retry behavior, with local identifiers and state transitions used to prevent duplicate logical completion in the tested workflow.

\subsection{Protected Transmission and Access}

Before a batch leaves the patient device, the client serializes the PGHD payload \(P\), generates an AES key \(k\) and nonce \(n\), and computes authenticated ciphertext

\begin{equation}
	c = \operatorname{AES\text{-}GCM.Enc}_{k,n}(P).
	\label{eq:encrypt}
\end{equation}

The client then computes \(h=\operatorname{SHA\text{-}256}(c)\) and signs the hash with the patient's ECDSA private key:

\begin{equation}
	\sigma = \operatorname{Sign}_{sk_p}(h).
	\label{eq:signature}
\end{equation}

Digital signatures provide integrity and source authentication when the corresponding public key is trusted \cite{NIST_FIPS186-5}.
The key and nonce are wrapped using the patient's PGHD PRE public key, producing encrypted key material and a PRE capsule.
The submission envelope contains the batch and patient identifiers, \(c\), \(h\), \(\sigma\), the wrapped key material, and the capsule.

Upon submission, PRE retrieves the patient's public key from IOTA, recomputes the ciphertext hash, and verifies the signature before any external write occurs.
Invalid envelopes are rejected.
For valid envelopes, PRE stores \(c\) in IPFS and records its CID, \(h\), validity status, and serialized cryptographic metadata in the patient's IOTA PGHD store.
Partitioning metadata into a separate \texttt{PatientPghdStore} per patient avoids forcing independent patients to mutate one global object during concurrent submissions.

To grant access, the patient creates an IOTA capability for a clinician and a re-encryption key for the corresponding PGHD purpose.
An authenticated clinician requests an entry through PRE.
PRE first verifies the token, capability owner, recipient, scope, validity interval, and revocation state.
Only then does it produce the re-encryption fragment required by the clinician.
The clinician retrieves the ciphertext, checks its CID and hash, verifies the patient's signature, reconstructs the symmetric key locally, and performs AES-GCM decryption.
Thus, plaintext is created only at an authorized endpoint.

Revocation prevents subsequent requests from obtaining the required re-encryption material.
It does not retroactively erase plaintext that an authorized clinician already accessed, which is an inherent limitation of data sharing rather than a property unique to PRE.

\subsection{Provenance and Invalidity Handling}

The encrypted payload uses three levels.
The batch level records the schema version, batch identifier, patient identifier, source device, covered time interval, and trigger reason.
The data-group level records the measurement type, source category, source package or device, recording method, statistics, thresholds, and anomaly count.
The data-point level stores the timestamp, value, unit, and anomaly markers.

This split limits public metadata while retaining clinically useful traceability for authorized readers.
IOTA provides the authoritative CID, access capability, and validity state, whereas detailed device and processing provenance remains encrypted in IPFS.
If PRE or the clinician detects a hash mismatch, invalid signature, incompatible schema, or inconsistent CID, the smart contract appends a new state transition that marks the entry invalid and records the failure reason.
The historical ledger transaction remains unchanged.

\subsection{Prototype Implementation}

The patient client was implemented in Kotlin using Jetpack Compose, Health Connect, Room, and WorkManager.
Rust shared libraries are accessed through the Java Native Interface for Umbral PRE operations and IOTA transaction construction, avoiding a second implementation of DecMed's cryptographic protocol.
The clinician client uses Tauri, SvelteKit, TypeScript, and Rust.
The PRE service uses Rust and Axum; Redis holds temporary access material; IPFS stores ciphertext; and IOTA Move smart contracts maintain accounts, capabilities, PGHD metadata, and validity state.
